\chapter{Introduction}

Inter-individual variability is ubiquitous in speech production and perception (see, e.g., \citealt{yuIndividualDifferencesLanguage2019} for a review of studies on individual variability in segmental processing). Prosody is no exception to this phenomenon. In fact, prosodic prominence, e.g., the employment of accent, is a system that lends itself especially readily to substantial variability. In the aptly titled journal paper ``Accent is predictable (if you're a mind reader)", \citet{bolingerAccentPredictableIf1972} expresses the idea that accent placement predominantly reflects a speaker's intent in structuring information.
Intent is naturally subjective and thus highly likely to give rise to variable productions across similar conditions. 
 
In West-Germanic languages, prominence encodes aspects of meaning, such as focus or information status \citep{buringSemanticsIntonationInformation2007, feryPitchAccentScaling2008, baumannCoreferenceLexicalGivenness2013, coleProsodyContextReview2015}. The most informative part of an utterance is understood as in \textit{focus}, other parts are in the \textit{background} \citep{buringSemanticsIntonationInformation2007}. Often the constituent in focus at the same time introduces \textit{new} information to the discourse, but this is not necessarily the case, so that information status can be regarded as an independent meaning dimension. Focus and information status can be subsumed under the general notion of information structure \citep{hallidayNotesTransitivityTheme1967}. Structuring information with regard to the preceding discourse is a basic communicative function that allows discourse participants to establish and manage common ground \citep{krifkaBasicNotionsInformation2008}. Parts of an utterance can be presented in focus or in the background, as contrastive or non-contrastive, and as new or given information \citep{buringSemanticsIntonationInformation2007}.

These meaning contrasts are encoded in the prominence \textit{relations} of prosodic constituents, which can be expressed by different (combinations of) prosodic parameters, such as changes in the fundamental frequency (F0), duration, and intensity \citep{fryDurationIntensityPhysical1955, sluijterSpectralBalanceAcoustic1996, kochanskiLoudnessPredictsProminence2005, baumannWhatMakesWord2018, roessigTracingPhoneticSpace2022}.
This potential redundancy allows individuals to rely on different aspects of the signal in prominence production and perception.
Relationality and redundancy can thus be regarded as key notions of prominence that enable and drive inter-individual variability.

In prosodic research, this variability has often been treated as noise and thus largely ignored \citep{arvanitiVariabilityOverlapCue2024}. With the advance of empirical methods, such as mixed-effects regression modeling, variability introduced by speakers or listeners is now usually accounted for in a random effects structure \citep[e.g.,][]{winterStatisticsLinguistsIntroduction2020, winterIndependenceGeneralizabilityLinguistics2021}. At least implicitly, inter-individual variability is thus acknowledged by most empirical studies that adhere to these standards, although generally only with the aim to control for this variability and thus clean up group-level effects.

Interest in individual variability from a theoretical point of view has also recently been on the rise. Promising lines of research address prosodic trade-off effects, cue weighting, and processing difficulties in the light of a lack of invariance of prosodic categories and form-meaning mapping \citep{kurumadaThinkingProbabilisticallyStudy2021, arvanitiVariabilityOverlapCue2024, severijnenYourVOORnaamNot2024}. Beyond these studies, explicit mentions of individual differences are also increasing in recent prosodic research, reflecting an effort towards greater empirical accuracy \citep[e.g.,][]{cangemiListenerSpecificPerceptionSpeakerSpecific2015, ouyangIndividualDifferencesProsodic2015, baumannWhatMakesWord2018, esteve-gibertEmpathyInfluencesHow2020}.

Nevertheless, systematic investigations into individual differences in prosody are still rare and thus the nature of this variability and mechanisms potentially underlying it are not yet well understood. Especially in the domain of prosody-meaning mapping, previous studies have concentrated on focus marking, neglecting other aspects of information structure. This book seeks to contribute to the emerging research concerned with individual variability in prosody production and perception. Specifically, we explore how and to what extent speakers and listeners vary in their encoding and decoding of prosodic prominence relations, both in terms of strength of encoding and types of parameters employed. Information status serves as the main prosodically encoded meaning aspect investigated here, but additionally, we consider syntactic factors that are known to interact with prosodic structure, namely grammatical role and sentence position. 

Furthermore, we take into account individuals' communicative skill as a potential influence on form-meaning mapping. Differences in communicative skill have previously been related to varying sensitivity to prosody \citep{bishopPhonologyPhoneticsSignalextrinsic2020, orricoInfluenceEmpathyAutisticlike2023} and pragmatic processing abilities \citep{nieuwlandIncrementalityPragmaticProcessing2010, zhaoAreScalarImplicatures2015} and are thus considered here as a potential explanation for inter-individual variability in prominence production and perception.

Empirical findings from two production studies and one perception study are presented and discussed in this book. These studies focus on eliciting speech and perception data from German native speakers in different information status conditions and syntactic structures. In the production experiments, participants produce highly scripted speech in reading tasks to allow for maximum control over the manipulations in the target utterances. To introduce some degree of naturalness, tasks are designed to be listener-directed and include an interactive component. In particular, the first production study investigates stories read by 32 speakers, containing target words that are either discourse new or accessible through the context. In the second production study, we analyze recordings of 20 speakers producing stories that contain target words varying in lexical and referential information status, grammatical role, and sentence position. In the perception study, we elicit prominence ratings on target words presented in context from 65 listeners to probe the effects of prosodic and non-prosodic effects on prominence perception. Here, we use recordings from the second production study as stimuli, including a variety of prosodic and syntactic structures as well as words representing different information status categories.

This book is firmly rooted in autosegmental-metrical theory \citep{laddIntonationalPhonology2008} and considers not only continuous but also categorical parameters associated with a prosodic hierarchy. This allows for direct implications for prosodic theory to be drawn. In terms of continuous measures, we rely on the recently proposed ProPer measures \citep{albertModelSonorityBased2023}, which are signal-based and derivable with minimal manual analysis steps. The same prosodic parameters are considered in all data analyses to allow for comparisons across studies.

In all three studies, we take an exploratory approach to the respective research questions. Exploratory data analysis as opposed to confirmatory analysis focuses on finding new patterns in the data rather than testing established hypotheses via rigorous hypothesis tests \citep{roettgerEmergentDataAnalysis2019}. There are two main reasons underlying the choice of an exploratory approach in this book. Firstly, individual differences in prosody-meaning mapping are not particularly well studied, so that formulating hypotheses with regard to the extent and nature of variability that we expect to observe is not a straightforward task. Secondly, some methods of statistical analysis that are used in this book are inherently exploratory, aiming to mine data for new patterns and not allowing for confirmatory hypothesis testing, e.g., random forest analysis and cluster analysis. Nevertheless, we will formulate expectations based on previous literature that will be tested against the data, but acknowledge that our results may well differ from these \citep{rubinExploratoryHypothesisTests2022}. Some aspects of these analyses are naturally more confirmatory than others, since more research is available to base assumptions on, especially with regard to general across-speaker analyses.

To accompany this exploratory approach, we conduct our statistical analyses in the Bayesian framework. While frequentist frameworks typically support a binary decision with regard to significance, the Bayesian framework allows for the quantification of uncertainty regarding the observed effects \citep{vasishthBayesianDataAnalysis2018}. Another practical advantage is that more complex models can be fitted in the Bayesian framework \citep{vasishthBayesianDataAnalysis2018}. Random slopes are often the first parameter to be dropped from frequentist mixed-effects regression models in cases of convergence issues \citep{barrRandomEffectsStructure2013}. However, these random slopes are crucial in the study of inter-individual variability, so that the Bayesian framework is deemed the best fit for the present investigation.

This book is structured as follows: Chapter \ref{ch:background prominence} outlines the concept of prominence underlying the current work and summarizes insights on prosodic structure, acoustic correlates of prominence, information structure, and interactions between prosody and lexical and syntactic factors, concluding with a discussion of the concept of redundancy in prominence marking and its implications for the emergence of individual variability. This chapter crucially motivates the parameters investigated in the production and perception experiments. Chapter \ref{ch:background variability} provides an overview of previous findings related to individual variability in prosody production and perception and discusses some potential causes, such as communicative skill, empathy, sociodemographic factors, and experience. Chapters \ref{ch:Production I} to \ref{ch: Perception} present the production and perception studies. In the first production study in Chapter \ref{ch:Production I}, we investigate individual variability in the prosodic encoding of information status. This is a first exploration of differences and similarities across speakers varying only one aspect of meaning, while at the same time comparing descriptive and more systematic analytic approaches to the data. The second production study presented in Chapter \ref{ch:Production II} expands on these findings by teasing apart the lexical and the referential level of information status and adding syntactic factors, namely grammatical role and sentence position. Communicative skill is also investigated as a potential factor modulating differences in prominence production across speakers. In Chapter \ref{ch: Perception}, we present a perception study that contributes to the understanding of individual variability with a prominence rating task. In particular, we explore how individuals differ with regard to their usage of prosodic and non-prosodic factors in decoding prominence relations. We again consider communicative skill as a potential underlying mechanism. Chapter \ref{ch: general discussion} brings together all findings in a general discussion and draws final conclusions on the main contributions of this book.